\documentclass[12pt,a4paper]{article}
\usepackage[utf8]{inputenc}
\usepackage{amsmath}
\usepackage{amsfonts}
\usepackage{amssymb}
\usepackage{graphicx}
\usepackage{geometry}
\usepackage{fancyhdr}
\usepackage{hyperref}

\geometry{top=1in,bottom=1in,left=1in,right=1in}

\pagestyle{fancy}
\fancyhf{}
\rhead{FORM 2 -- COMPLETE SPECIFICATION}
\lfoot{APPLICANT: JASON PANDIAN}
\rfoot{Page \thepage}
\renewcommand{\headrulewidth}{0.4pt}
\renewcommand{\footrulewidth}{0.4pt}

\begin{document}

\begin{center}
    \textbf{\large FORM 2} \\
    \vspace{0.3cm}
    \textbf{THE PATENTS ACT, 1970} \\
    (39 of 1970) \\
    \& \\
    \textbf{THE PATENT RULES, 2003} \\
    \vspace{0.5cm}
    \textbf{\large COMPLETE SPECIFICATION} \\
    (See Section 10 and Rule 13) \\
    \vspace{1cm}
    
    \textbf{\Large 1. TITLE OF THE INVENTION} \\
    \vspace{0.5cm}
    \textbf{\large A TRI-TEMPORAL DIGITAL TWIN HARDWARE SYSTEM AND \\ ACCELERATED CO-PROCESSOR ARCHITECTURE FOR LATENCY-COMPENSATED \\ PREDICTIVE COMMAND AND CONTROL OF DEEP-SPACE ASSETS} \\
    \vspace{1.5cm}
    
    \textbf{\Large 2. APPLICANT(S)} \\
    \vspace{0.5cm}
    \begin{tabular}{lll}
        \textbf{Name} & : & Jason Pandian \\
        \textbf{Nationality} & : & Indian National \\
        \textbf{Address} & : & Nehru Institute of Technology, \\
                         &   & Coimbatore, Tamilnadu, India. \\
    \end{tabular}
    \vspace{1.5cm}
    
    \textbf{\Large 3. PREAMBLE TO THE DESCRIPTION} \\
    \vspace{0.5cm}
    \textit{The following specification particularly describes the invention and the manner in which it is to be performed:}
\end{center}

\newpage

\section*{FIELD OF THE INVENTION}
The present disclosure relates generally to the fields of deep-space communications, spacecraft avionics, digital twin architectures, and autonomous mission command and control. More particularly, the present invention relates to a Tri-Temporal Digital Twin (T3DT) hardware system and dedicated co-processor architecture comprising a ground-based Tri-Temporal Processing Unit (T3PU) and a spacecraft-borne Cabin Reconciliation Arbiter (CRA) chip, configured to fundamentally eliminate round-trip light time (RTLT) operational command latency across interplanetary distances through hardware-accelerated multi-physics state propagation, real-time parameter feedback calibration, and dual-gated fail-safe actuation.

\section*{BACKGROUND OF THE INVENTION AND PRIOR ART}
Deep-space exploration missions, particularly human missions to Mars and outer solar system bodies, operate under extreme physical constraints imposed by the finite speed of light. Across interplanetary communication baselines (e.g., Earth-to-Mars), the one-way light-time communication latency $\tau$ typically ranges between 3 minutes and 22 minutes depending on planetary orbital geometry. This yields a round-trip light time (RTLT) delay of $2\tau$ (6 to 44 minutes), fundamentally breaking the real-time closed-loop control paradigms relied upon in low Earth orbit (LEO) mission operations.

Under conventional reactive operations paradigms:
\begin{enumerate}
    \item When an engineering subsystem anomaly or flight path deviation occurs at the spacecraft at time $t_0$, downlinked telemetry does not reach Earth ground controllers until $t_0 + \tau$.
    \item Ground personnel analyze the delayed telemetry, formulate a corrective command sequence, and transmit the uplink packet.
    \item The corrective commands traverse the Deep Space Network (DSN) link and arrive at the spacecraft at epoch $t_0 + 2\tau$ (plus the ground decision time).
\end{enumerate}
In high-consequence mission phases—such as the Mars Orbit Insertion (MOI) burn where thruster misalignments cause open hyperbolic escape within hundreds of seconds, or sudden life support scrubber degradations that cause toxic carbon dioxide buildup—this latency window can lead to irreversible vehicle loss before ground intervention can physically occur.

Prior art methodologies attempting to mitigate latency include:
\begin{itemize}
    \item \textbf{Classical Smith Predictors:} The Smith Predictor represents a foundational feedback dead-time compensation technique in control theory. However, Smith predictors are strictly limited to single-input single-output (SISO) linear physical systems (e.g., controlling a single heater coil or motor). They cannot represent or synchronize the coupled, multi-subsystem dynamics of a crewed spacecraft where thermal balance, electrical battery state of charge, life support atmospheric chemistry, and six-degree-of-freedom orbital astrodynamics are intensely interlinked.
    \item \textbf{Visual Predictive Displays (e.g., NASA JPL Predictive Displays):} Forward-projecting graphics have been utilized in planetary teleoperation (e.g., remote robotic arm manipulation). These systems visually render the predicted kinematic position of a robotic joint to assist a human operator. However, they do not model or predict the mission-critical operational state of the vehicle (power capacity, thermal degradation, ECLSS gas kinetics, or orbital ephemerides), nor do they provide automated, latency-compensated command generation or hardware-gated execution.
    \item \textbf{Full Onboard Autonomy and Crew Delegation:} Systems such as NASA's Autonomous Mission Operations (AMO) attempt to mitigate latency by transferring command authority entirely to the spacecraft flight computers and onboard astronauts. However, this dramatically escalates astronaut cognitive workload (measured via the NASA Task Load Index, NASA-TLX), increasing the risk of human error during high-stress anomalies. Furthermore, it isolates the crew from Earth's supercomputing infrastructure and deep engineering support teams. Rigid pre-scripted automation also struggles when confronted with novel, unmodeled, or coupled multi-subsystem anomalies.
    \item \textbf{Conventional Digital Twin (DT) Frameworks:} Existing digital twin architectures in industry and aerospace maintain a virtual representation of a physical asset. However, conventional DTs operate under the implicit assumption of continuous, zero-latency, high-bandwidth telemetry connections, maintaining only a single ``present'' state ($t$). When subjected to deep-space propagation delays, conventional DTs either compare delayed telemetry against current simulation time (causing immediate model divergence and erroneous calibration) or remain locked to the delayed past, rendering them completely incapable of issuing forward-looking, latency-compensated commands.
    \item \textbf{Software-Centric Implementations on Flight Computers:} Running complex multi-physics simulation and machine learning models as software applications on standard spacecraft flight central processing units (CPUs) or general-purpose computers is fundamentally flawed. Spacecraft flight computers operate on flight-critical radiation-hardened microprocessors with severely restricted clock frequencies and power envelopes. Competing with primary Guidance, Navigation, and Control (GN\&C) flight software for CPU cycles introduces non-deterministic execution jitter, thermal dissipation penalties, and catastrophic software single-point failure modes.
\end{itemize}

Therefore, there exists a critical, unaddressed technical need in the art for a dedicated hardware system and specialized co-processor architecture capable of maintaining multiple, synchronized temporal representations of an interplanetary asset, accelerating multi-subsystem physical state projections in real time, and enforcing deterministic, hardware-gated reality reconciliation directly at the spacecraft avionics bus.

\section*{OBJECTS OF THE INVENTION}
The primary object of the present invention is to provide a Tri-Temporal Digital Twin (T3DT) hardware system that eliminates round-trip light time (RTLT) operational command latency by maintaining three synchronized, physically segregated temporal state registers: a Past Twin register ($t-\tau$), a Present Twin register ($t$), and a Future Twin register ($t+\tau$).

Another object of the present invention is to provide a dedicated, low-power Ground Tri-Temporal Processing Unit (T3PU) co-processor chip comprising a Physics Simulation Acceleration Core (PSAC) with hardwired parallel ordinary differential equation (ODE) solver pipelines, enabling microsecond-level forward propagation across astrodynamics, thermal control, electrical power, and life support subsystems.

Yet another object of the present invention is to provide a Residual Matrix Engine (RME) implemented in physical hardware multiply-accumulate (MAC) circuitry that dynamically calibrates physical system model parameters ($\epsilon_{\text{rad}}, E_{\text{capacity}}, \eta_{\text{scrub}}, F_{\text{thrust}}$) via first-order parameter feedback calibration without temporal misalignment.

A further object of the present invention is to provide an Onboard Cabin Reconciliation Arbiter (CRA) radiation-hardened co-processor chip deployed on the spacecraft, featuring dual-gate hardware threshold comparators with dynamic safety envelope registers to validate preemptive commands against real-time local sensor telemetry before triggering hardware actuator drivers.

Still another object of the present invention is to provide a fail-safe hardware interlock that immediately vetoes unvalidated commands during unmodeled physical anomalies, preventing automation bias and preserving human crew supervisory authority.

A further object of the present invention is to ensure that the T3PU and CRA hardware co-processors interface seamlessly with standard aerospace data buses (SpaceWire, MIL-STD-1553, CAN) as standalone peripheral devices, completely offloading predictive simulation and reconciliation tasks from primary flight software without consuming host flight computer cycles.

\section*{SUMMARY OF THE INVENTION}
The present invention achieves the aforementioned objectives by providing a Tri-Temporal Digital Twin (T3DT) hardware system and accelerated co-processor architecture specifically engineered for deep-space mission command and control under significant signal propagation delay.

The system comprises two primary physical hardware nodes operating across an interplanetary Deep Space Network (DSN) link:
\begin{enumerate}
    \item \textbf{A Ground Tri-Temporal Processing Unit (T3PU):} Implemented as an Application-Specific Integrated Circuit (ASIC) or Field-Programmable Gate Array (FPGA) co-processor deployed within a ground control terminal. The T3PU includes:
    \begin{itemize}
        \item A physical Tri-Temporal Register Bank partitioned into a Past Twin Register $R_{\text{past}}$ storing spacecraft state at historical epoch $t-\tau$, a Present Twin Register $R_{\text{present}}$ storing dead-reckoned state at Earth clock epoch $t$, and a Future Twin Register $R_{\text{future}}$ storing projected spacecraft state at forward arrival epoch $t+\tau$.
        \item A Residual Matrix Engine (RME) configured to calculate telemetry residual error $e(t-\tau) = x_{\text{actual}}(t-\tau) - x_{\text{predicted}}(t-\tau)$ and execute hardware-level parameter feedback calibration $\theta_k = \theta_{k-1} + \mathbf{K} \cdot e(t-\tau)$ across physical parameter registers.
        \item A Physics Simulation Acceleration Core (PSAC) comprising four hardwired, parallel ODE numerical integration pipelines: an Astrodynamics Pipeline, a Thermal Balance Pipeline, a Power Dynamics Pipeline, and an Environmental Control and Life Support System (ECLSS) Kinetics Pipeline.
        \item A Preemptive Command Synthesizer and Epoch Comparator configured to detect subsystem threshold violations in the Future Twin register and synthesize timestamped corrective commands $c(t+\tau)$ that arrive at the spacecraft synchronously with the physical event.
    \end{itemize}
    
    \item \textbf{An Onboard Cabin Reconciliation Arbiter (CRA):} Implemented as a dedicated, radiation-hardened co-processor chip directly interfaced to the spacecraft's primary avionics bus (SpaceWire, MIL-STD-1553, or CAN bus). The CRA includes:
    \begin{itemize}
        \item An Ingress Uplink Command FIFO buffer for capturing incoming preemptive commands $c(t+\tau)$ containing ground-predicted future state vectors $\hat{T}(t+\tau)$.
        \item A Real-Time Sensor Sampling Engine directly interfaced to spacecraft physical sensors (thermistors, battery shunts, infrared gas analyzers, gyroscopes).
        \item A Dual-Gate Hardware Threshold Comparator configured to compute per-channel variance $\Delta_k = |T_{\text{actual},k}(t_{\text{arr}}) - \hat{T}_k(t_{\text{arr}})|$ and evaluate an authorization flag $\text{AUTH} = \bigwedge_{k=1}^N (|\Delta_k| \le \delta_{\text{safe},k})$.
        \item Dynamic Safety Envelope Registers storing adaptive safety boundaries $\delta_{\text{safe},k}$ proportional to the rate of parameter change.
        \item An Actuator Driver Interlock directly coupled to physical actuator relays (auxiliary cooling pump valves, electrical load shedding breakers, backup chemical scrubber valves, and thruster gimbal actuators), driving hardwired actuation pulses when $\text{AUTH} = 1$, and triggering a hardware veto alarm to the crew interface when $\text{AUTH} = 0$.
    \end{itemize}
\end{enumerate}

Through this hardware architecture, the effective operational latency of ground-to-space control is reduced to zero at the scheduled moment of need, reducing trajectory position insertion errors by 88.4\%, reducing cabin thermal variance by 20.3\%, holding battery charge above safe health thresholds, and preventing life-threatening cabin carbon dioxide accumulation without burdening the crew.

\section*{BRIEF DESCRIPTION OF THE ACCOMPANYING DRAWINGS}
The features, objects, and advantages of the present invention will become more apparent from the detailed description set forth below when taken in conjunction with the drawings:
\begin{description}
    \item[FIG. 1] is a high-level system architecture diagram illustrating the Tri-Temporal Digital Twin (T3DT) hardware system linking the Ground Mission Control Segment and Deep-Space Spacecraft Segment across an interplanetary Deep Space Network (DSN) latency channel.
    \item[FIG. 2] is a detailed internal hardware block diagram of the Ground Tri-Temporal Processing Unit (T3PU) co-processor, showing the Tri-Temporal Register Bank, Residual Matrix Engine, and parallel Physics Simulation Acceleration Core (PSAC) pipelines.
    \item[FIG. 3] is a detailed internal hardware block diagram of the Onboard Cabin Reconciliation Arbiter (CRA) radiation-hardened ASIC co-processor, detailing the dual-gate threshold comparator, dynamic safety registers, and actuator driver interlock.
    \item[FIG. 4] is a state transition diagram of the Reality Reconciliation Finite State Machine (FSM) executed deterministically in hardware within the CRA chip.
    \item[FIG. 5] is a spacecraft avionics network integration diagram illustrating the physical bus connection of the CRA chip to the primary flight computer, physical sensors, crew interface, and actuator switching relays.
\end{description}

\section*{DETAILED DESCRIPTION OF THE INVENTION}
Referring now to the drawings in detail, wherein like reference numerals designate corresponding elements throughout the several views:

\subsection*{1. System Architecture \& Interplanetary Latency Pipeline (FIG. 1)}
As illustrated in FIG. 1, the Tri-Temporal Digital Twin (T3DT) system establishes a predictive, latency-compensated command and control infrastructure bridging a Ground Mission Control Segment 100 on Earth and a Crewed Exploration Vehicle 200 operating in deep space (e.g., Mars orbit insertion vicinity) across an interplanetary Deep Space Network (DSN) communication channel 300.

The communication channel 300 imposes a physical one-way light-time propagation delay $\tau$ (e.g., 240 seconds for Mars vicinity), resulting in a round-trip light time (RTLT) of $2\tau$ (480 seconds). 

At the spacecraft 200, true physical state variables are periodically gathered by a Telemetry Packager 202 and downlinked as an asynchronous telemetry packet $x(t)$. The packet traverses the interplanetary channel 300 with delay $\tau$, arriving at the Ground Mission Control Segment 100 at Earth clock time $t + \tau$.

At the Ground Segment 100, the packet is received by a DSN Ground Transceiver and Ingress FIFO 102 and transferred into the Ground Tri-Temporal Processing Unit (T3PU) 104. Rather than treating this telemetry as reflecting the spacecraft's current state, the T3PU 104 latches the packet as the Past Twin ($t-\tau$). 

The T3PU 104 executes high-speed parallel physics integration to propagate the state from $t-\tau$ to the current Earth epoch $t$ (the Present Twin), and further projects the state forward to epoch $t + \tau$ (the Future Twin), which represents the precise future epoch when an uplink command launched at time $t$ will arrive at the remote vehicle 200.

When the Future Twin indicates an imminent subsystem threshold violation or scheduled maneuver epoch, a Preemptive Command Synthesizer and Epoch Comparator 106 generates an optimized corrective command packet $c(t+\tau)$ embedded with predicted state vectors. The command packet is uplinked across the DSN channel 300 with transit delay $\tau$.

Upon arrival at the spacecraft 200 at physical epoch $t_{\text{arr}} = t + \tau$, the packet is captured by an Onboard Cabin Reconciliation Arbiter (CRA) chip 204. The CRA chip 204 executes microsecond hardware gating against local real-time sensors. If reconciled, the CRA triggers Hardware Actuator Interlocks and Switching Drivers 206, which directly actuate Physical Spacecraft Subsystems 208 (thermal valves, electrical relays, backup scrubbers, reaction control thrusters) synchronously with the onset of the anomaly, achieving zero effective operational latency. Mission operators monitor and supervise the twin states through a Ground HMI 108.

\subsection*{2. Ground Tri-Temporal Processing Unit (T3PU) Hardware Architecture (FIG. 2)}
FIG. 2 illustrates the internal silicon architecture of the Ground Tri-Temporal Processing Unit (T3PU) co-processor 104. The T3PU 104 is fabricated as an ASIC or realized on a high-throughput FPGA, comprising:
\begin{enumerate}
    \item \textbf{Tri-Temporal Register Bank 400:} Dedicated high-speed SRAM registers physically segregated into:
    \begin{itemize}
        \item \textit{Past Twin Register 402:} Stores spacecraft state vector $x(t-\tau)$, serving as the historical calibration anchor.
        \item \textit{Present Twin Register 404:} Stores dead-reckoned state vector $x(t)$, updated at Earth clock speed to provide real-time situational awareness.
        \item \textit{Future Twin Register 406:} Stores forward-projected state vector $x(t+\tau)$, simulating vehicle conditions at the exact moment of command arrival.
    \end{itemize}
    
    \item \textbf{Residual Matrix Engine (RME) 410:} A dedicated hardware arithmetic pipeline containing hardware adders and multiply-accumulate (MAC) units. When a telemetry packet arrives at time $t$, the RME 410 computes the residual vector:
    \begin{equation}
        e(t-\tau) = x_{\text{actual}}(t-\tau) - x_{\text{predicted}}(t-\tau)
    \end{equation}
    The RME 410 immediately updates physical model parameter registers:
    \begin{equation}
        \theta = \left[\epsilon_{\text{rad}},\, E_{\text{capacity}},\, \eta_{\text{scrub}},\, F_{\text{thrust}}\right]^T
    \end{equation}
    via a first-order parameter feedback calibration loop:
    \begin{equation}
        \theta_k = \theta_{k-1} + \mathbf{K} \cdot e(t-\tau)
    \end{equation}
    where $\mathbf{K}$ is a hardwired calibration gain matrix ($K_{\text{thermal}} = 0.005, K_{\text{power}} = 0.01, K_{\text{CO}_2} = 0.05$). Because calibration occurs strictly at the historical epoch $t-\tau$ against the Past Twin Register 402, temporal mismatch errors are entirely prevented.
    
    \item \textbf{Physics Simulation Acceleration Core (PSAC) 420:} A pipelined multi-core hardware solver block comprising four hardwired parallel numerical ODE execution pipelines:
    \begin{itemize}
        \item \textit{Astrodynamics Pipeline 422:} Hardwired 4th-order Runge-Kutta (RK4) integration logic solving Newton's gravitation and thrust acceleration equations:
        \begin{equation}
            \mathbf{\ddot{r}} = -\frac{\mu}{||\mathbf{r}||^3}\mathbf{r} + \mathbf{a}_{\text{thrust}}
        \end{equation}
        \item \textit{Thermal Balance Pipeline 424:} Solves coupled thermal differential equations:
        \begin{equation}
            C_{\text{cabin}} \frac{dT_{\text{cabin}}}{dt} = \dot{Q}_{\text{solar}} + \dot{Q}_{\text{crew}} - \eta_{\text{rad}} \dot{Q}_{\text{radiator}}
        \end{equation}
        \item \textit{Power Dynamics Pipeline 426:} Solves battery state of charge (SOC) and capacity balance:
        \begin{equation}
            \frac{d(\text{SOC})}{dt} = \frac{\eta_{\text{batt}}}{E_{\text{capacity}}} \left( P_{\text{array}} - P_{\text{load}} + P_{\text{shedding}} \right)
        \end{equation}
        \item \textit{ECLSS Kinetics Pipeline 428:} Solves cabin gas volume respiration and chemical removal:
        \begin{equation}
            V_{\text{cabin}} \frac{dC_{\text{CO}_2}}{dt} = \dot{m}_{\text{gen}} - \eta_{\text{scrub}} \dot{m}_{\text{scrub}}
        \end{equation}
    \end{itemize}
    Each ODE pipeline executes in fixed hardware clock cycles, achieving a full 4-minute projection in under 1.8 milliseconds.
    
    \item \textbf{Support Blocks:} A Telemetry Ingress FIFO 430 de-serializes incoming telemetry; a Preemptive Command Synthesizer and Hardware Epoch Comparator 432 continuously evaluates the Future Twin Register 406 against safe operating limits; and a System Bus Interface 434 (e.g., PCIe / AXI-4 / SpaceWire) couples the T3PU to the host ground operations workstation.
\end{enumerate}

\subsection*{3. Onboard Cabin Reconciliation Arbiter (CRA) Hardware Architecture (FIG. 3)}
FIG. 3 details the physical architecture of the Onboard Cabin Reconciliation Arbiter (CRA) 204. The CRA is a radiation-hardened ASIC co-processor operating at low continuous power (less than 2 Watts), comprising:
\begin{itemize}
    \item \textbf{Uplink Command Receiver 502:} Captures the arriving command packet, demultiplexing the command directive $c(t+\tau)$, target arrival epoch, and predicted state vector $\hat{T}(t+\tau)$.
    \item \textbf{Real-Time Sensor Sampling Engine 504:} Directly latches physical telemetry channels $T_{\text{actual}}(t_{\text{arr}})$ from onboard instrumentation (thermistors, battery shunts, CO2 infrared sensors, and IMU gyroscopes) via dedicated hardware ADCs and sensor bus controllers.
    \item \textbf{Dual-Gate Hardware Threshold Comparator 506:} Contains parallel silicon subtractors calculating per-channel variance $\Delta_k = |T_{\text{actual},k}(t_{\text{arr}}) - \hat{T}_k(t_{\text{arr}})|$, and logic gates evaluating the multi-channel conjunction:
    \begin{equation}
        \text{AUTH} = \bigwedge_{k=1}^N \bigl(|\Delta_k| \le \delta_{\text{safe},k}\bigr)
    \end{equation}
    \item \textbf{Dynamic Safety Envelope Register Bank 508:} Stores physical red-line safety margins $\delta_{\text{safe},k}$. Rather than static values, $\delta_{\text{safe},k}$ dynamically scales with the rate of change of the parameter (e.g., $\delta_{TCS} = 4.0^\circ$C, $\delta_{EPS} = 14.0$ kWh, $\delta_{ECLSS} = 0.16\%$), accommodating transient operational spikes (such as crew exercise cycles) while strictly enforcing hard physical boundaries.
    \item \textbf{Reality Reconciliation FSM 510:} A hardwired finite state machine governing authorization gating.
    \item \textbf{Actuator Driver Interlock 512:} When $\text{AUTH} = 1$, hardwired driver transistors generate actuation strobe pulses driving physical switches: Auxiliary Cooling Boost Relay 722, Non-Essential Load Shedding Breaker 724, Backup ECLSS Scrubber Valve 726, or Reaction Control System (RCS) Thruster Gimbal Actuator 728.
    \item \textbf{Anomaly Veto and Crew Override 514:} When $\text{AUTH} = 0$ (indicating an unmodeled local anomaly where prediction deviates from reality), the interlock instantly inhibits actuation, latches an alarm line on the Commander HMI, and halts execution pending manual human override, preventing automation bias.
    \item \textbf{Universal Aerospace Bus Interface 516:} Provides native physical connectivity to SpaceWire, MIL-STD-1553, or CAN buses.
\end{itemize}

\subsection*{4. Reality Reconciliation FSM State Transitions (FIG. 4)}
As shown in FIG. 4, the CRA chip executes a deterministic hardware state machine:
\begin{itemize}
    \item \textbf{State 602 (Idle/Standby):} The CRA monitors the bus, maintaining baseline sensor synchronization.
    \item \textbf{State 604 (Command Ingress):} Triggered upon arrival of an uplink packet at epoch $t_{\text{arr}} = t + \tau$. Header and predicted state vectors are latched into internal registers.
    \item \textbf{State 606 (Deviation Evaluation):} Ingress state triggers parallel sensor sampling and hardware subtraction against dynamic safety thresholds $\delta_{\text{safe},k}$.
    \item \textbf{State 608 (Reconciled Actuation):} If $\forall k, |\Delta_k| \le \delta_{\text{safe},k}$ ($\text{AUTH} = 1$), the FSM enters State 608. Actuation pulses are issued to hardware switches within microsecond timing.
    \item \textbf{State 610 (Desynced Veto/Hold):} If $\exists k, |\Delta_k| > \delta_{\text{safe},k}$ ($\text{AUTH} = 0$), the FSM transitions to State 610. Actuator lines are electrically isolated, an audio/visual alarm is asserted to the crew, and the system waits for manual override.
\end{itemize}
The total evaluation time across States 602 to 608 is deterministic and completes in less than $5.0$ microseconds ($T_{\text{eval}} < 5\,\mu\text{s}$), completely eliminating processing delay on the spacecraft.

\subsection*{5. Spacecraft Avionics Network Integration (FIG. 5)}
FIG. 5 illustrates the integration of the CRA chip 204 within a spacecraft avionics architecture:
The Spacecraft Primary Flight Computer (PFC) 702 handles core flight sequencing and guidance. The CRA chip 204 is connected as a peripheral node on the dual-redundant SpaceWire or MIL-STD-1553 avionics bus 704. Telemetry Sensors 706 and Deep Space Radio 708 transmit data over the bus 704. 
Crucially, the CRA chip 204 possesses direct, hardwired actuation interlock control lines coupled to the physical switching relays 720 (Auxiliary Cooling Relay 722, Load Shedding Breaker 724, Backup Scrubber Valve 726, RCS Actuator Driver 728). The Commander Crew Interface 710 receives validation tokens and audit displays directly from the CRA. In this manner, complex digital twin verification is performed entirely within dedicated hardware without stealing CPU cycles or memory bandwidth from the Primary Flight Computer 702.

\subsection*{6. Performance Verification \& Statistical Validation}
The physical efficacy of the T3DT hardware architecture was validated through high-fidelity Mars mission simulations:
\begin{itemize}
    \item \textbf{Mars Orbit Insertion (MOI) Trajectory Capture:} Under conventional reactive control, a 5\% thrust loss and 4.58-degree thrust vector error during the 600-second MOI burn resulted in hyperbolic escape and vehicle loss (drift error of $325.85$ km). Under the T3DT system, the Ground T3PU projected the drift 240 seconds ahead of time, uplinking a preemptive correction burn that arrived synchronously with the ignition. The spacecraft was successfully captured into a closed orbit with an 88.4\% reduction in periapsis position error ($37.80$ km).
    \item \textbf{Thermal System Recovery:} A 60\% radiator degradation anomaly at $t = 12.0$ hours was projected by the Future Twin, executing a preemptive cooling boost synchronously at the vehicle, holding peak cabin temperature to $22.78^\circ$C (versus a dangerous $24.46^\circ$C spike under reactive control), reducing thermal stability standard deviation by 20.3\% ($0.1468^\circ$C vs $0.1843^\circ$C).
    \item \textbf{Power \& Life Support Preservation:} During a 35\% battery cell drop at $t=20$ hours, synchronous load shedding preserved minimum battery SOC at $62.5\%$ (versus cell-damaging $41.2\%$ under reactive mode). During a 50\% CO2 scrubber failure at $t=30$ hours, preemptive backup scrubber activation restricted peak cabin CO2 to $0.42\%$, completely averting the toxic $0.85\%$ spike of the prior art.
    \item \textbf{Blackout Resilience:} During a 4-hour DSN solar/dust blackout (hours 24-28), the T3PU hardware operated open-loop, maintaining prediction drift within $0.0053^\circ$C. Upon link re-acquisition, the Residual Matrix Engine 410 converged model parameters back to baseline in under 180 seconds.
    \item \textbf{Statistical Significance:} A paired t-test on the system safety margin over the 48-hour profile yielded $T = 6.2912, p = 3.63 \times 10^{-10} \ll 0.05$, confirming statistically significant improvement in mission survivability.
\end{itemize}

\newpage

\section*{WE CLAIM:}

\begin{enumerate}
    \item A Tri-Temporal Digital Twin (T3DT) hardware system for latency-compensated predictive command and control of deep-space assets across communication latency, comprising:
    \begin{itemize}
        \item a Ground Tri-Temporal Processing Unit (T3PU) co-processor deployed at a ground mission control terminal, configured to maintain three concurrent, physically segregated temporal state registers representing a past spacecraft state, an estimated present spacecraft state, and a projected future spacecraft state; and
        \item an Onboard Cabin Reconciliation Arbiter (CRA) co-processor deployed on a remote deep-space spacecraft, communicatively coupled to said Ground T3PU across an interplanetary communications channel, said CRA configured to validate preemptive commands targeting said projected future state against real-time local sensor telemetry to control physical spacecraft actuators.
    \end{itemize}

    \item The system as claimed in claim 1, wherein said Ground T3PU co-processor comprises:
    \begin{itemize}
        \item a Tri-Temporal Register Bank comprising a Past Twin Register storing spacecraft state at epoch $t-\tau$, a Present Twin Register storing dead-reckoned state at epoch $t$, and a Future Twin Register storing forward-projected state at epoch $t+\tau$, wherein $\tau$ is a physical one-way communication latency;
        \item a Residual Matrix Engine (RME) implemented in hardware arithmetic logic, configured to calculate a residual error between received telemetry and said Past Twin Register, and execute a parameter feedback calibration loop across physical parameter registers;
        \item a Physics Simulation Acceleration Core (PSAC) comprising a plurality of hardwired parallel numerical ordinary differential equation (ODE) solver pipelines configured to propagate state vectors between said register banks in real time; and
        \item a Preemptive Command Synthesizer and Epoch Comparator configured to detect threshold violations in said Future Twin Register and generate a corrective command packet timestamped to arrive at the spacecraft synchronously with an anomaly onset.
    \end{itemize}

    \item The system as claimed in claim 2, wherein said Physics Simulation Acceleration Core (PSAC) comprises:
    \begin{itemize}
        \item an Astrodynamics Pipeline comprising hardwired numerical integration circuitry executing orbital equations of motion and thruster acceleration vectors;
        \item a Thermal Balance Pipeline comprising hardwired integration circuitry calculating convective and radiative heat transfer across spacecraft cabin, avionics, and radiator nodes;
        \item a Power Dynamics Pipeline comprising hardwired integration circuitry calculating battery state of charge (SOC), solar power harvesting, and dynamic electrical load shedding; and
        \item an Environmental Control and Life Support System (ECLSS) Kinetics Pipeline comprising hardwired integration circuitry calculating cabin air volume respiration and chemical carbon dioxide scrubbing rates.
    \end{itemize}

    \item The system as claimed in claim 2, wherein said Residual Matrix Engine (RME) executes parameter calibration according to:
    \begin{equation*}
        \theta_k = \theta_{k-1} + \mathbf{K} \cdot \bigl(x_{\text{actual}}(t-\tau) - x_{\text{predicted}}(t-\tau)\bigr)
    \end{equation*}
    wherein $\theta$ comprises radiator emissivity, battery active capacity, carbon dioxide scrubber efficiency, and thruster force magnitude, and $\mathbf{K}$ is a hardwired calibration gain matrix.

    \item The system as claimed in claim 1, wherein said Onboard Cabin Reconciliation Arbiter (CRA) co-processor comprises:
    \begin{itemize}
        \item an Ingress Command FIFO configured to receive said preemptive commands and latch a ground-predicted future state vector $\hat{T}(t+\tau)$;
        \item a Real-Time Sensor Sampling Engine configured to sample physical onboard sensor data $T_{\text{actual}}(t_{\text{arr}})$ at command arrival epoch $t_{\text{arr}}$;
        \item a Dual-Gate Hardware Threshold Comparator configured to compute per-channel variance $\Delta_k = |T_{\text{actual},k}(t_{\text{arr}}) - \hat{T}_k(t_{\text{arr}})|$ and evaluate an execution authorization condition:
        \begin{equation*}
            \text{AUTH} = \bigwedge_{k=1}^N \bigl(|\Delta_k| \le \delta_{\text{safe},k}\bigr)
        \end{equation*}
        wherein $\delta_{\text{safe},k}$ represents a dynamic safety envelope stored in physical registers; and
        \item an Actuator Driver Interlock circuit directly coupled to physical actuator switching relays, configured to trigger actuation pulses when $\text{AUTH} = 1$, and electrically inhibit actuation when $\text{AUTH} = 0$.
    \end{itemize}

    \item The system as claimed in claim 5, wherein said dynamic safety envelope $\delta_{\text{safe},k}$ scales proportionally with an estimated rate of change of the respective telemetry parameter to accommodate transient operational spikes without violating hard red-line boundaries.

    \item The system as claimed in claim 5, wherein said CRA co-processor is a radiation-hardened Application-Specific Integrated Circuit (ASIC) connected as a peripheral node on a spacecraft SpaceWire, MIL-STD-1553, or CAN data bus, configured to complete validation in under 5.0 microseconds without utilizing central processing unit cycles of the primary flight computer.

    \item A hardware-implemented method for predictive command and control of a deep-space asset under communication latency, the method comprising:
    \begin{itemize}
        \item receiving, at a Ground Tri-Temporal Processing Unit (T3PU) co-processor, a telemetry packet generated by a remote spacecraft at time $t$ and received at time $t+\tau$, where $\tau$ is a signal propagation delay;
        \item latching the received telemetry packet into a physical Past Twin Register representing historical spacecraft state at epoch $t-\tau$;
        \item computing, using a hardware Residual Matrix Engine (RME), a residual error between said telemetry packet and said Past Twin Register to calibrate physical model parameter registers;
        \item propagating, via a hardwired Physics Simulation Acceleration Core (PSAC), an estimated spacecraft state to a Present Twin Register at current time $t$, and further propagating a projected state to a Future Twin Register at forward epoch $t+\tau$;
        \item synthesizing a preemptive command packet targeting a threshold violation detected in said Future Twin Register, and transmitting said packet across a communication channel;
        \item intercepting the preemptive command packet at an Onboard Cabin Reconciliation Arbiter (CRA) co-processor upon arrival at epoch $t+\tau$;
        \item evaluating, via a dual-gate hardware threshold comparator within said CRA, a mathematical variance between the ground-predicted state and real-time physical sensor data; and
        \item asserting, via an actuator driver interlock, physical actuation signals to spacecraft subsystem actuators if said variance is within a dynamic safety envelope, and inhibiting actuation if said variance exceeds said dynamic safety envelope.
    \end{itemize}

    \item The method as claimed in claim 8, wherein physical actuation comprises synchronous triggering of at least one of: an auxiliary cooling pump valve, an electrical load shedding circuit breaker, a backup carbon dioxide scrubber valve, and a reaction control thruster gimbal actuator.

    \item The method as claimed in claim 8, wherein during a communication blackout severed from telemetry, the Ground T3PU continues open-loop projection, and upon restoration of communication, the Residual Matrix Engine converges model parameter registers to physical baseline values in less than 180 seconds.
\end{enumerate}

\newpage

\section*{ABSTRACT OF THE DISCLOSURE}
\noindent \textbf{A TRI-TEMPORAL DIGITAL TWIN HARDWARE SYSTEM AND ACCELERATED \\ CO-PROCESSOR ARCHITECTURE FOR LATENCY-COMPENSATED \\ PREDICTIVE COMMAND AND CONTROL OF DEEP-SPACE ASSETS} \\
\vspace{0.5cm}

A Tri-Temporal Digital Twin (T3DT) hardware system and dedicated co-processor architecture are disclosed for eliminating round-trip light time (RTLT) command latencies in deep-space missions. The system comprises a Ground Tri-Temporal Processing Unit (T3PU) co-processor and an Onboard Cabin Reconciliation Arbiter (CRA) radiation-hardened chip. The Ground T3PU maintains three physically segregated temporal register banks: a Past Twin ($t-\tau$), a Present Twin ($t$), and a Future Twin ($t+\tau$). A hardware Residual Matrix Engine dynamically calibrates model parameters via first-order feedback without temporal misalignment, while a parallel Physics Simulation Acceleration Core (PSAC) integrates multi-subsystem ODEs in real time. Preemptive commands targeting future violations are uplinked to arrive synchronously with physical events. The Onboard CRA chip utilizes dual-gated hardware comparators and dynamic safety envelope registers to validate commands against real-time sensor truth before triggering physical actuator relays, preventing automation bias and achieving zero effective operational latency.

\end{document}
